\begin{table}[H]
\centering
\scriptsize
\renewcommand{\arraystretch}{1.35}
\begin{tabularx}{\textwidth}{|c|X|X|X|X|X|}
\hline
\textbf{Sprint} & \textbf{Objectif principal} & \textbf{Frontend / Mock-first} & \textbf{Backend / Intégration} & \textbf{Livrables} & \textbf{Critères de sortie} \\
\hline
\textbf{Sprint 1\\Fondations UX + Mock} &
Valider les parcours clés de l'application avec des données simulées et figer les contrats techniques. &
- Initialiser React Native bare + TypeScript.
- Mettre en place navigation, thème, composants UI réutilisables.
- Créer les écrans : onboarding, login, register, forgot password, home feed, post detail, create post, profile.
- Implémenter un store local et un dossier \texttt{services/mock}.
- Simuler utilisateurs, posts, commentaires, followers, conversations.
- Ajouter upload image simulé, feed chronologique simulé, likes/commentaires simulés. &
- Initialiser FastAPI, PostgreSQL, Redis, Alembic, Docker Compose.
- Créer modèles de base : User, Follow, Post, PostMedia, Comment, Like.
- Implémenter auth JWT, refresh token, hash bcrypt.
- Préparer structure modulaire backend et conventions d'API.
- Documenter les DTO v1 pour alignement frontend/backend. &
- App mobile navigable de bout en bout en mode mock.
- Backend démarrable localement avec auth prête.
- Contrats API versionnés.
- Maquette fonctionnelle présentable. &
- Un utilisateur peut parcourir toute l'app sans backend réel.
- Les types frontend sont figés.
- Le backend permet déjà register/login/me.
- L'équipe valide les parcours UX avant branchement réel. \\
\hline
\textbf{Sprint 2\\Social + Intégration Progressive API} &
Remplacer les mocks des fonctionnalités sociales par le backend réel sans casser l'UX validée. &
- Connecter login/register/refresh au backend.
- Remplacer progressivement le feed mock par les vraies données.
- Brancher création de post, affichage profil, follow/unfollow, likes, commentaires.
- Garder des fallbacks mock pour les écrans non encore branchés.
- Ajouter écran de vote ``Tenue A vs Tenue B'' côté UI. &
- Finaliser endpoints : auth, profile, upload média, posts, feed, comments, likes, follows, polls.
- Intégrer stockage MinIO/S3.
- Ajouter pagination feed.
- Contrôles d'accès et validation des uploads.
- Préparer conversations directes et structure de messages. &
- Parcours social principal connecté au backend réel.
- Création de posts avec images fonctionnelle.
- Vote communautaire disponible.
- Documentation API sociale mise à jour. &
- Un utilisateur réel peut s'inscrire, publier, suivre, liker, commenter.
- Le feed se charge depuis PostgreSQL.
- Les uploads image fonctionnent sur environnement local/staging.
- Les mocks restants sont clairement isolés. \\
\hline
\textbf{Sprint 3\\Temps Réel + WebRTC Progressif} &
Livrer le temps réel en plusieurs niveaux : chat stable, appel audio, puis vidéo. &
- Implémenter UI de conversations 1:1.
- Brancher Socket.IO côté mobile pour messages, présence, accusés de lecture.
- Ajouter partage d'un post dans le chat.
- Créer écran d'appel avec états : ringing, connecting, in-call, ended.
- Intégrer d'abord l'appel audio, puis activer la vidéo et le switch caméra. &
- Mettre en place Socket.IO backend.
- Gérer conversations directes, messages texte/image/post partagé.
- Ajouter présence online/offline et statuts sent/delivered/read.
- Implémenter signalisation WebRTC via Socket.IO.
- Déployer Coturn et tester ICE/STUN/TURN.
- Journaliser les erreurs d'appel et timeouts. &
- Chat temps réel 1:1 complet.
- Appel audio 1:1 stable.
- Appel vidéo 1:1 fonctionnel sur cas nominal.
- Rapport de tests réseau réels. &
- Deux utilisateurs peuvent chatter en temps réel.
- Les messages image et post partagé fonctionnent.
- L'appel audio passe sur Wi-Fi et 4G.
- La vidéo fonctionne au moins sur scénario nominal avec TURN de secours. \\
\hline
\textbf{Sprint 4\\IA + Finition + Démo} &
Remplacer les derniers mocks critiques par l'IA réelle et préparer une version de démonstration stable. &
- Ajouter bouton ``Help Me Choose'' sur post/photo.
- Créer écran résultats IA avec 3 cards visuelles.
- Permettre le partage d'une recommandation IA dans le chat.
- Ajouter loaders, retry, états d'erreur et UX de fallback.
- Finaliser polish visuel et cohérence navigation. &
- Implémenter endpoint \texttt{/ai/help-me-choose}.
- Ajouter service IA avec mode mock/dev et mode OpenAI/staging.
- Construire prompt pipeline avec historique, occasion, style.
- Mesurer latence et protéger par timeout.
- Finaliser sécurité, rate limiting, logs, configuration staging HTTPS.
- Préparer script seed de démonstration. &
- Assistant IA intégré.
- Démo MVP stable de bout en bout.
- Environnement staging présentable.
- Jeu de données de démonstration prêt. &
- Un utilisateur peut demander un avis IA et recevoir 3 suggestions.
- La réponse IA reste dans la cible de temps prévue ou bascule vers un fallback propre.
- L'application est démontrable sans intervention manuelle technique.
- Tous les parcours critiques ont été testés avant soutenance. \\
\hline
\end{tabularx}
\caption{Roadmap DressMe -- approche progressive mock-first puis intégration backend}
\end{table}
